\documentclass[11pt]{article}
\usepackage[T1]{fontenc}
\usepackage{textcomp}
\usepackage[margin=0.75in]{geometry}
\usepackage{amsmath,amssymb,amsthm}
\usepackage{array,booktabs}
\usepackage{hyperref}
\setlength{\parindent}{0pt}
\newtheorem{theorem}{Theorem}
\theoremstyle{definition}
\newtheorem{definition}{Definition}
\title{Flow Proof Artifact: right-triangle-acute-angles-sum}
\author{Generated by \texttt{flow doc proof}}
\date{Flow Proof Book}
\begin{document}
\maketitle
The two acute angles in a right triangle sum to one right angle.\\[0.5em]
\textbf{Source.} euclid — Elements, Book I, Proposition 32\\[0.5em]
\bigskip
\noindent{\large \textbf{Derived fact 1.} \textit{angle A plus angle B equals one right angle when angle C is right} \hfill the Euclidean plane \cdot right triangle \cdot acute angles sum to one right angle}\par
\smallskip\noindent\textit{Coordinate: the Euclidean plane \cdot right triangle \cdot acute angles sum to one right angle}\par
\smallskip\noindent\textit{Source: euclid}\par
\smallskip\noindent\textit{Built on: an angle inscribed in a semicircle is a right angle, the three interior angles of a triangle sum to two right angles}\par
\medskip
\noindent\textbf{Goal.} angle A plus angle B equals one right angle when angle C is right\par
\noindent$\displaystyle \angle A + \angle B = 90^\circ$\par
\medskip
\noindent
\begin{tabular}{@{} >{\raggedright\arraybackslash}p{0.06\textwidth} >{\raggedright\arraybackslash}p{0.47\textwidth} >{\raggedleft\arraybackslash}p{0.41\textwidth} @{}}
\textbf{\#} & \textbf{Proof} & \textbf{Mathematics} \\
\midrule
\textbf{1.} & We prove that acute angles sum to one right angle for right triangles in the Euclidean plane. & \\[0.45em]
\textbf{2.} & We invoke the derived fact governing circles in the Euclidean plane: an angle inscribed in a semicircle is a right angle. & \\[0.45em]
\textbf{3.} & We invoke the derived fact governing triangles in the Euclidean plane: the three interior angles of a triangle sum to two right angles. & \\[0.45em]
\textbf{4.} & From step 2 and step 3, this implies angle A plus angle B equals one right angle. Hence proven. & $\displaystyle \angle A + \angle B = 90^\circ$ \\[0.45em]
\bottomrule
\end{tabular}
\label{thm:Geometry-right triangle-acute_angles_sum_to_one_right_angle}
\par\medskip\hrule
\end{document}
